\chapter{Fundamentação Teórica} \label{Cap:Fund_Teo}


\section{Mercado Financeiro Brasileiro}

%\subsection{Hipótese de Mercados Eficientes e Hipótese de Mercados Adaptativos}
%\subsection{Formação de Preços e Assimetria Informacional}

\subsection{Índices de Ações da B3}

\subsubsection{Ibovespa (IBOV)}

O Ibovespa é o principal indicador de desempenho das ações negociadas na B3 e reúne as empresas mais importantes do mercado de capitais brasileiro. Foi criado em 1968 e de lá para cá consolidou-se como referência para investidores ao redor do mundo.

Reavaliado a cada quatro meses, o índice é composto pelas ações e units de companhias listadas na B3 que correspondem a cerca de 85\% do número de negócios e do volume financeiro do nosso mercado de capitais.

\subsubsection{Índice Dividendos (IDIV)}

Objetivo de muitos dos investidores, os dividendos também têm sua participação medida em índice. O objetivo do IDIV é indicar o desempenho médio dos ativos que se destacam na distribuição de remuneração aos investidores, seja sob a forma de dividendos ou juros sobre o capital próprio (JCP).

\subsubsection{Índice Small Caps (SMLL)}

O objetivo do SMLL é ser o indicador do desempenho médio de uma carteira composta por ações de empresas de menor capitalização. Ele é composto pelas ações e units que estão fora da lista do Ibovespa, que representa 85\% do valor de mercado de todas as empresas na B3.

% \section{Processamento de Linguagem Natural (PLN)}
% \subsection{Conceitos Fundamentais de PLN}
% \subsection{Modelos de Linguagem Baseados em Transformers}
% \subsubsection{Arquitetura BERT}
% \subsubsection{Modelos para Português: BERTimbau}
% \subsubsection{Modelos Específicos de Domínio: FinBERT-PT-BR}
% \subsection{Análise de Sentimento}
% \subsubsection{Classificação de Polaridade (Positivo, Negativo, Neutro)}
% \subsubsection{Análise de Sentimento em Textos Financeiros}


% \section{Séries Temporais Financeiras}
% \subsection{Características de Séries Temporais Financeiras}
% \subsubsection{Não Estacionariedade}
% \subsubsection{Volatilidade e Heterocedasticidade}
% \subsubsection{Autocorrelação}
% \subsection{Modelos Tradicionais}
% \subsubsection{Modelos ARIMA}
% \subsubsection{Modelos GARCH}
% \subsection{Modelos de Deep Learning para Séries Temporais}
% \subsubsection{Redes Neurais Recorrentes (RNN)}
% \subsubsection{Long Short-Term Memory (LSTM)}
% \subsubsection{Gated Recurrent Unit (GRU)}
% \subsubsection{Transformers para Séries Temporais}


% \section{Modelos Híbridos: Integração de PLN e Séries Temporais}
% \subsection{Abordagens de Fusão de Features}
% \subsubsection{Fusão Precoce (\textit{Early Fusion})}
% \subsubsection{Fusão Tardia (\textit{Late Fusion})}
% \subsubsection{Fusão Híbrida}
% \subsection{Arquiteturas Multimodais para Predição Financeira}


% \section{Big Data e Mídias Sociais em Finanças}
% \subsection{Características de Big Data (Volume, Velocidade, Variedade)}
% \subsection{X/Twitter como Fonte de Dados Financeiros}
% \subsection{Desafios de Processamento de Dados Textuais Massivos}


% \section{Trabalhos Relacionados}
% \subsection{Análise de Sentimento no Mercado Internacional}
% \subsection{Análise de Sentimento no Mercado Brasileiro}
% \subsection{Modelos Híbridos de Predição de Preços}